% Created 2019-06-08 Sat 10:57
% Intended LaTeX compiler: pdflatex
\documentclass[11pt]{article}
\usepackage[utf8]{inputenc}
\usepackage[T1]{fontenc}
\usepackage{graphicx}
\usepackage{grffile}
\usepackage{longtable}
\usepackage{wrapfig}
\usepackage{rotating}
\usepackage[normalem]{ulem}
\usepackage{amsmath}
\usepackage{textcomp}
\usepackage{amssymb}
\usepackage{capt-of}
\usepackage{hyperref}
\author{emsenn}
\date{\today}
\title{30 Days of World Building, Summer 2019}
\hypersetup{
 pdfauthor={emsenn},
 pdftitle={30 Days of World Building, Summer 2019},
 pdfkeywords={},
 pdfsubject={},
 pdfcreator={Emacs 26.2 (Org mode 9.2.3)},
 pdflang={English}}
\begin{document}

\maketitle
\section*{Introduction}
\label{sec:org6a18707}
This collection of writing exercises was written in response to the \emph{30 Days of World
Building} challenge created by \emph{Stephanie Bryant} from June 4th until
July 4th.

This collection of writing exercises was created by me, emsenn, in the United States. To the extent possible under law, I have waived all copyright and related or neighboring rights to this document. This document was created for the benefit of the public. If you're viewing it on a remote server, you're encouraged to download your own copy. To learn how you can support me, see \href{https://emsenn.net/donate}{https://emsenn.net/donate}.
\section*{30 Days of World Building, Summer 2019}
\label{sec:org3d2e47d}
\startcontents[level-1]
\printcontents[level-1]{}{0}{\setcounter{tocdepth}{2}}
\subsection*{Day 1}
\label{sec:org617eb65}
The exercise for this first day is this:

\begin{quote}
Get out a map or go to an international website like National
Geographic. Look every- where. Antarctica. Saudi Arabia. The
rainforests of Brazil. The rainforests of Central California. Look
at how the different climates behave and appear.

The first fifteen minute exercise is to write down all the different
climates you can think of-- if you need to just say a city name, do
it. Sometimes "Seattle" is more evocative than "northern damp
temperate climate." Write these names down in a list.

Then, go through that list and write one or two words that describe
how that climate, either the word itself or the way the place itself
may have made you feel, if you've been there before. Try to stick to
abstract adjectives; emotional words, if you can, but nouns are also
okay.
\end{quote}

Simple enough - though I think they might mean "biome" where they say
"climate." And I'm not sure all of what I've listed are biomes,
"officially."

\begin{center}
\begin{tabular}{ll}
Climate & Description\\
\hline
alpine mountains & divisive\\
badlands & sanguine, ironic\\
bayou & calm, saturated\\
chaparral & endurance, commitment\\
coral reef & childishness, riddles\\
equatorial archipelago & allegorical, precise\\
fen & playful\\
floodlands & folkart, communion\\
freshwater lake & plain-spoken\\
grasslands & eternal\\
holler & sneakiness\\
lowlands & historic, homely\\
mangrove swamp & consuming, decaying\\
marsh & honor\\
meadows & betrayal, beauty\\
moors & treason\\
ocean & introspection, philosophy\\
old growth forest & treacherous, ebulent\\
piedmont & prosperity, equality\\
pine forest & snide, caustic\\
rainforest & aggrevating, inspiring\\
redwood forest & weird\\
river delta & rigid, eternal\\
sandy desert & barren, spiritual\\
savanna & funny, excited\\
sea & opulence, reservation\\
shield & hardworking\\
steppes & stagnation\\
taiga & enthralling, catatonic\\
thicket & surprise\\
tidewater & creativity\\
tundra & challenging, introspective\\
\end{tabular}
\end{center}

\subsection*{Day 2}
\label{sec:orgf7d8828}
Before I get to today's exercise, I want to express a concern, of
sorts, I have about being able accurately complete this challenge. I
read through future days, and it seems that the challenge comes from
the assumption that you have a vague idea of a story, and not a clear
idea of what sort of world to have.

I am coming at this from rather the opposite perspective: I have a
notion of the world, but am looking for more stories to tell in it.
This may mean that, going forward, I'll have to adjust how I respond
to certain exercises.

The exercise for this second day is this:

\begin{quote}
What role do you anticipate weather playing in your story? Do you have
a lot of travel that you want to complicate with bad storms? Are you
going to snow in your mighty heroes? Will there be a mighty battle,
determined by sudden flooding? You actually don't need to know right
now. Your exerise for today is to jot down ten plot devices that
relate to weather, and what you think they do to the story (for
example: a snow-in can turn the mood very claustrophobic... or very
intimate).
\end{quote}

Okay, ten plot devices related to weather, and what they can do to a
story.

\begin{center}
\begin{tabular}{ll}
Device & Effect\\
\hline
calm at sea & Slowly building desperation. It's hard to know how long a calm will last when it starts, and so it becomes a point of anxiety, as it stretches on.\\
flash flood & Sudden emergency. You've got to make choices about what to rescue, where to go, right now.\\
fog & Obfuscation, opportunity for invisible things. Unseen murders.\\
weeks of rain & Discomfort, duldrums, break from routine. Canceled parties, missed opportunity.\\
heat wave & Concern for the future. Will it kill the crops? You've got time to plan.\\
thunderstorm & Natural fears. They're just... scary.\\
dust storm & Confusion, isolation, lasting destruction.\\
spring shower & Opportunity for romance, as two seek shelter from the sudden downpour.\\
very sunny & Noticing things you might otherwise not. The glint of a bracelet as a thief removes it from an aristocrat's wrist in the market.\\
snowstorm & Holiday special! Or other community-based things.\\
\end{tabular}
\end{center}

I feel I could dive into any of these to turn them into an \emph{adventure}
or story set within \emph{Teraum}. One reminds me of an idea I had
recently, for an \emph{adventure} set on a transoceanic cargo ship upon
which the \emph{party} has bought passage. Instead of scenes being defined
by area, they'd be defined by what the main challenge each month is,
and one of those months would be the calm month, in which the ship
is becalmed, and the \emph{party} has to figure out how to remedy that
situation, by getting the ship to move or stretching the supplies.

Just for example.
\subsection*{Day 3}
\label{sec:orged77828}
The exercise for the third day is this:

\begin{quote}
Read over your list from Day 1 and then turn to a blank piece of
paper. Close your eyes and think about what kind of feeling you like
to have when you write or read. Write down four words that fit into
that feeling: two adjectives, a verb, and a noun. Now return to the
page with your list of climates and emotions. Do any of them match up?
If they do, you have your climate. If not, try to find closest-match
words.

If you spend 10 solid minutes thinking about this and still can't
decide, pick two climates that express moods you like. You can make up
your mind later, and you can even build your world with both climates
containing equally probable sites for your story.
\end{quote}

What kind of story do I want to tell? One where people are desperate
and innovative; one where reactive decision-making ruins plans; one
where even the heroes are covered in mud, but probably not blood:

\begin{itemize}
\item desperate
\item innovative
\item react
\item mud
\end{itemize}

So looking at the table from the first day's exercise, a few biomes
stick out to me:

\begin{itemize}
\item I say "tidewater," is a setting for creativity - that's reasonably
close to innovative. I also say "rainforests" are inspiring.
\item "Thickets" are associated with surprise for me, that's close to
react.
\item "Freshwater lake" has plain-spoken, and floodlands has folkart, both
of which are about the closest I can find in my list to mud.
\end{itemize}

So, where should the story be set, following this logic? Sounds like
someplace that's normally subject to the tides, with lots of little
bits of forest to hideaway in.

There's one part of my world that is very much like that: \emph{Arathel
County}, a coastal region north of humanity's largest city. A mix of
wetlands, the whole region shifts with the tide, and again with spring
flooding.

\subsection*{Day 4}
\label{sec:org0e755d3}
The excercise for this fourth day is:

\begin{quote}
For 15 minutes, jot down some of the Really Big Land Features you want
in your story and just think "what if that were made by...." Write
down a couple of causes for those features and scars, and stick it all
into your notebook. Which scars are slow-force scars (like plate
tectonics), and which are fast-forces (anything that takes less than
10,000 years is medium-to-fast in geological scales).
\end{quote}

In yesterday's prompt, I was led to decide that \emph{Arathel County} would
be where I'm setting the stories and world-building these exercises
are bringing me too.

That might make this challenge a bit tricky, since it's catered around
\emph{world}-building, not \emph{small region}-building. In talking about this
hallenge with a couple friends doing it, they also expressed some
difficulty - the challenge seems to be drifting between "develop a
world" and "develop the area of the world your story will be set in."

Anyway. The biggest land feature in the region is the region itself.
It's on the eastern (and only known) coast of the \emph{Old World}, north
of the \emph{Green River}, in a part of the world known as the \emph{Green
Delta}.

The \emph{Old World}, as written so far, has four major \emph{tectonic plates}.
The \emph{Green Delta} and the \emph{Central Plains} to its east exist on one,
which is subducting the \emph{Optic Ocean}, the ocean to the region's
west.

To the southeast, the \emph{Central Plains} are obducting over the \emph{Unseen
Sea}, a massive desert which now exists a couple thousand feet under
sea level; the tectonic border has formed a massive cliff known as the
\emph{Falling Coast}.

To the east, the \emph{Central Plains} and \emph{Unseen Sea} meet with the
\emph{Worldkeepers}, converging to form that mountain range.

So, four plates: under the \emph{Green Delta} and \emph{Central Plains}, under
the \emph{Optic Ocean}, under the \emph{Unseen Sea}, and under the
\emph{Worldkeepers}.

That brings us back to \emph{Arathel County}, which, despite being part of
the \emph{Green Delta}, is not on the same plate. The region was formerly a
coral reef and archipelago of sandy shoals, in the \emph{Optic Ocean}.

As the plates shifted, the region got obducted - I'm not sure if
that's the proper word anymore, but "moved to be ontop of," - by the
shifting of the \emph{Green Delta} plate.

An aside: it's interesting not having names for these plates because
the people in my setting wouldn't be aware of them. Also, I should
point out, while I'm using "plate tectonics" as the term, and the
world does exist as a crust of plates, the plates' cause for movement
is not like Earth's. Underneath the world is, basically, a zoo, called
the \emph{Worldkeep}, for all sorts of magic creatures. It's quite old by
now and parts are collapsing in, and this is causing the world above
to form plates and for them to move.

Anyway, \emph{Arathel County} was a coral reef, but shifting of the plates
caused it to end up above the shoreline. The coral got filled in with
sediment over many years, and the sandy islands got worn down to dirt.

This has led to a region that's mostly gravel and sandy, with complex
and often inter-connected cave regions. There's lots of ferns that
grow in the region, and moss, and those sorts of things.

There's a river that runs through the river, the \emph{Kaiper River},
though "runs" is an over-ambitious description. Fed exclusively by
rainfall and springs in \emph{Arathel County}, the river is low-volume and
full of silt, slowly oozing toward the coast.

That's about it for notable geographic features in the region. The
town of \emph{Arathel} is built on the coast of the \emph{Kaiper River}, and
over time the area directly adjacent to the road leading to the town
from the rest of the \emph{Green Delta} has been cultivated into small
scrappy farms.
\stopcontents[level-1]
\section*{Supplements}
\label{sec:org63a4c1b}
\subsection*{Index}
\label{sec:org75d290d}
\begin{description}
\item[{30 Days of World Building}] A series of prompts for building a fictional world. See See \href{http://www.web-writer.net/fantasy/days/}{http://www.web-writer.net/fantasy/days/} or \#30DaysOfWorldBuilding.
\item[{Bryant, Stephanie}] Creator of the \emph{30 Days of World Building} challenge, and an Agile and SCRUM coach. See \href{http://www.mortaine.com/}{http://www.mortaine.com/}
\item[{plate tectonics}] A theory describing the movement of rock segments under a plates surface leading to the formation of geographic features. See \href{https://en.wikipedia.org/wiki/Plate\_tectonics}{https://en.wikipedia.org/wiki/Plate\_tectonics}.
\item[{Teraum}] A tragically funny fantasy setting. See See \href{https://emsenn.net/teraum/}{https://emsenn.net/teraum/}  or \#Teraum.
\item[{Teraum, Arathel}] A sleepy, worn-down town.
\item[{Teraum, Arathel County}] A swampy and sandy county.
\item[{Teraum, Central Plains}] A large agricultural region, full of towns and castles.
\item[{Teraum, Falling Coast}] A massive cliff dividing a vast desert from the rest of the world.
\item[{Teraum, Green Delta}] A hospitable and well-settled region of the world.
\item[{Teraum, Green River}] A slow wide river, cutting through the world.
\item[{Teraum, Kaiper River}] A silty excuse for a river.
\item[{Teraum, Old World}] The large continent where most of history happened.
\item[{Teraum, Optic Ocean}] A large desert separating one part of the world from the other.
\item[{Teraum, Unseen Sea}] A desert so large it was first mistaken for a seabed.
\item[{Teraum, Worldkeepers}] A range of large mountains.
\item[{Teraum, Worldkeep}] A vast complex housing an unbelievable variety of magic lifeforms.
\end{description}
\subsection*{README}
\label{sec:org7a75f02}
\subsubsection*{30 Days of World Building, Summer 2019}
\label{sec:orgd923f1d}
This collection of writing exercises was written in response to the \emph{30 Days of World
Building} challenge created by \emph{Stephanie Bryant} from June 4th until
July 4th.

This collection of writing exercises was created by me, emsenn, in the United States. To the extent possible under law, I have waived all copyright and related or neighboring rights to this document. This document was created for the benefit of the public. If you're viewing it on a remote server, you're encouraged to download your own copy. To learn how you can support me, see \href{https://emsenn.net/donate}{https://emsenn.net/donate}.
\end{document}
